\begin{abstract}
跨机构数据流通正在从单纯的数据传输转向“数据可用不可见、过程可追溯、责任可界定”的协作模式。BLoP 将区块链、可信计算与隐私保护计算结合，为数据流通中的可信节点治理、异常监测和 PBFT 风格链上确认提供了完整的全节点可信审计框架。面向高频审计和轻量化部署需求，本文以 BLoP 的可信节点集合和链上确认思想为基础，提出一种轻量化升级选项 TrustFlowAudit。该方法将复杂校验、证据存储和异常恢复放在链下执行，将数据承诺、审计摘要、恢复状态和节点信誉变化轻量写入联盟链；同时利用可信评分选择小规模审计委员会，减少每次审计均由全部可信节点参与确认带来的通信与链上记录开销。方法层面，TrustFlowAudit 结合 Merkle 批量存证、哈希链证据组织、纠删码恢复触发和 PBFT 风格联盟链确认，形成“分片存储、抽样审计、异常发现、恢复处理、链上记录、信誉更新”的闭环。本文暂不展开实验结果，重点给出问题定义、系统模型、方法流程和安全开销分析，为后续复现实验和小论文定量评估奠定基础。
\end{abstract}

\begin{IEEEkeywords}
数据流通，区块链，分布式数据校验，数据审计，可信委员会，PDP，PoR，纠删码，PBFT，Merkle 树
\end{IEEEkeywords}
